\documentclass[10pt,a4paper]{amsart}
\usepackage[margin=19mm]{geometry}
\usepackage{amsmath,amssymb,mathtools}
\usepackage[scr=rsfs,cal=euler]{mathalfa}
\usepackage{xfrac}
\usepackage[unicode,hidelinks,bookmarks=false]{hyperref}
\newcommand{\ie}{i.e.\ }
\newcommand{\cf}{cf.\ }
\newcommand{\resp}{respectively\ }
\newcommand{\Subsection}{\subsection}
\providecommand{\nobreakdash}{\nobreak\hbox{-}}
\setlength{\emergencystretch}{2em}
\multlinegap0pt
\usepackage{amssymb}
\usepackage[shortlabels]{enumitem}
\usepackage{xcolor}
\newtheorem{lemma}{Lemma}
\title{Radial Nodal Dirichlet Solutions of Singular Elliptic Equations: Global Branches, Endpoint Asymptotics, and Morse Indices}
\author{Wenjing Chen\ \ and Zexi Wang}
\date{Source: arXiv:2609.15283v1 (CC BY 4.0)}
\begin{document}
\maketitle

\noindent\emph{Source and licence.} Radial Nodal Dirichlet Solutions of Singular Elliptic Equations: Global Branches, Endpoint Asymptotics, and Morse Indices. Version arXiv:2609.15283v1 (CC BY 4.0), distributed by arXiv under the Creative Commons Attribution 4.0 International licence, http://creativecommons.org/licenses/by/4.0/, as recorded in the arXiv licence metadata for this exact version and shown on the abstract page https://arxiv.org/abs/2609.15283v1. Full licence notice: \texttt{licenses/arxiv-2609.15283-CC-BY-4.0.txt}.\par\smallskip

\noindent\textit{Editorial scope.} Counted unit: the article's lemma lem:sublinear-bessel-limit (source0.tex lines 1964--1981). The article's complete original proof of it is delivered with it (source0.tex lines 1983--2133, one \texttt{proof} environment of 151 lines). No mathematics is written, repaired or re-derived: the statement and the proof are the article's own lines, in the article's own notation, and no retained line is substituted or re-flowed.  Retained uncounted as the source-local context the proof needs: lines 362--366; lines 443--448; lines 1450--1478; lines 1738--1741.  Nothing was omitted from any retained span, so the package carries no omission record.  No retained line contains a citation, so the wrapper carries no bibliography and no bibliography entry.  No retained line contains a figure environment, an image-inclusion command or drawing code, so no image is delivered and the package declares no asset dependency.  Editorial formatting only, in addition: an AMS \texttt{amsart} preamble that restates the article's own declarations and macros used by the retained lines, namely the environments \texttt{lemma} on separate counters, together with the standard packages those lines need.  The retained environments print in this document as Lemma 1 in order of appearance, whereas the article prints the same items with the numbering of its own full document.  The complete native source of the article as distributed in the arXiv e-print for this exact version is bundled unchanged as \texttt{sources/210406/source0.tex}.
\begin{equation}\label{eq:main-bessel}
\Phi''+\frac{n-1}{s}\Phi'+\Phi=0,
\qquad \Phi(0)=1,
\qquad \Phi'(0)=0,
\end{equation}

\begin{equation}\label{eq:sub-constant-intro}
 \mathfrak c_{j,q,n}
 :=\frac{\displaystyle\int_0^{\rho_j}
 s^{n-1}|\Phi_n(s)|^q\,ds}
 {\rho_j^{n-1}\Phi_n'^2(\rho_j)}>0.
\end{equation}

\begin{lemma}\label{lem:log-bessel}
As $\alpha\to\infty$, equivalently $A\to\infty$, one has
$w_A\to\Phi_n$ in $C^1([0,S])$ for every $S>0$. Let $\Psi_n$ be the
unique regular integral solution of
\begin{equation}\label{eq:log-second-profile}
 \Psi''+\frac{n-1}{s}\Psi'+\Psi
 =-\Phi_n\log\Phi_n^2,
 \qquad \Psi(0)=\Psi'(0)=0.
\end{equation}
Then
\begin{equation}\label{eq:log-second-profile-convergence}
 A(w_A-\Phi_n)\longrightarrow\Psi_n
 \quad\hbox{in }C^1([0,S])
\end{equation}
for every $S>0$. Consequently, if $\rho_j$ is the $j$-th positive zero
of $\Phi_n$, then
\begin{equation}\label{eq:zeroasymptotic}
 \sqrt A\,z_j(\alpha)
 =\rho_j+\frac{c_{j,n}}{A}+o(A^{-1}),
 \qquad
 c_{j,n}:=-\frac{\displaystyle\int_0^{\rho_j}
 s^{n-1}\Phi_n^2(s)\log\Phi_n^2(s)\,ds}
 {\rho_j^{n-1}\Phi_n'^2(\rho_j)}>0.
\end{equation}
In particular,
\begin{equation}\label{eq:z-zero}
 \lim_{\alpha\to\infty}z_j(\alpha)=0.
\end{equation}
\end{lemma}

\begin{equation}
 z_{j}'(\alpha)<0.
 \label{eq:zero-curve-derivative}
\end{equation}

\begin{lemma}
\label{lem:sublinear-bessel-limit}
For every fixed $j\ge1$, the zero curve \(z_j\), defined on
\((\alpha_{j-1},\infty)\), satisfies
\begin{equation}
 z_j(\alpha)\downarrow\rho_j
 \quad\text{as }\alpha\to\infty.
 \label{eq:sublinear-right-limit-general}
\end{equation}
More precisely,
\begin{equation}
 \lim_{\alpha\to\infty}\alpha^{2-q}\bigl(z_j(\alpha)-\rho_j\bigr)
 = \mathfrak c_{j,q,n},
 \label{eq:sublinear-first-order-zero}
\end{equation}
where $\mathfrak c_{j,q,n}$ is defined in
\eqref{eq:sub-constant-intro}.
\end{lemma}

\begin{proof}
For $\alpha>1$, set
\[
 w_\alpha(r):=\frac{u(r,\alpha)}{\alpha},
 \qquad
 \varepsilon_\alpha:=\alpha^{q-2}.
\]
Since $1<q<2$, we have $\varepsilon_\alpha\rightarrow0$ as $\alpha\to\infty$.
If the shooting solution reaches a first double zero, we use its zero extension beyond that point. Thus $w_\alpha$ satisfies
\begin{equation}
 \begin{cases}
 w_\alpha''+\dfrac{n-1}{r}w_\alpha'+w_\alpha
 -\varepsilon_\alpha|w_\alpha|^{q-2}w_\alpha=0,\qquad  r>0,\\[4pt]
 w_\alpha(0)=1,\qquad w_\alpha'(0)=0.
 \end{cases}
 \label{eq:sublinear-scaled}
\end{equation}
  The energy inequality gives
$|u(r,\alpha)|<\alpha$ before the first double zero. Hence, with the zero extension,
\begin{equation}
|w_\alpha(r)|\leq1
\quad\text{for all }r\geq0.
\label{eq:upper1}
\end{equation}
The Volterra representation of \eqref{eq:sublinear-scaled} is
\begin{align}
 w_\alpha(r)=1-
 \int_0^r t^{1-n}\int_0^t s^{n-1}
 \Bigl(w_\alpha(s)-\varepsilon_\alpha
 |w_\alpha(s)|^{q-2}w_\alpha(s)\Bigr)\,ds\,dt.
 \label{eq:sublinear-volterra}
\end{align}
In particular,
\begin{equation*}
w_\alpha'(r)=
-r^{1-n}\int_0^r s^{n-1}
\Bigl(
w_\alpha(s)
-\varepsilon_\alpha
|w_\alpha(s)|^{q-2}w_\alpha(s)
\Bigr)ds.
\end{equation*}
Since {$|w_\alpha|\leq1$} and $\varepsilon_\alpha\rightarrow0$,
\begin{equation*}
\big|w_\alpha-\varepsilon_\alpha
|w_\alpha|^{q-2}w_\alpha\big|\leq |w_\alpha|+\varepsilon_\alpha|w_\alpha|^{q-1}\leq 1+\varepsilon_\alpha\leq2,
\end{equation*}
and therefore,
\begin{equation*}
|w_\alpha'(r)|\leq
2r^{1-n}\int_0^r s^{n-1}ds\leq \frac{2}{n}r.
\end{equation*}
{For $r>0$, the equation also gives
\[
 |w_\alpha''(r)|
 \le \frac{n-1}{r}|w_\alpha'(r)|
 +\bigl|w_\alpha(r)-\varepsilon_\alpha
 |w_\alpha(r)|^{q-2}w_\alpha(r)\bigr|
 \le \frac{2(n-1)}{n}+2.
\]
The same estimate at $r=0$ follows from the regular integral
formulation. Hence $\{w_\alpha'\}$ is equi-Lipschitz on every fixed
compact interval.}
Hence, for every
$L>0$, the families $\{w_\alpha\}$ and $\{w_\alpha'\}$ are uniformly
bounded and equicontinuous on $[0,L]$. The Arzel\`{a}--Ascoli theorem
therefore yields relative compactness of $\{w_\alpha\}$ in
$C^1([0,L])$. Let $\tilde w$ be the $C^1$-limit of
a subsequence. Passing to the limit in \eqref{eq:sublinear-volterra}, and using
$\varepsilon_{\alpha}\to0$, we find that $\tilde w$ solves \eqref{eq:main-bessel}. By the uniqueness of the regular solution, $\tilde w=\Phi_n$. Since every convergent subsequence has the same limit, the whole family satisfies
\begin{equation}\label{eq:y-to-Phi}
w_\alpha\rightarrow\Phi_n\quad \text{in }C^1_{\rm loc}([0,\infty)){\quad\text{as }\alpha\to\infty}.
\end{equation}
By the $C^1_{\mathrm{loc}}$ convergence in \eqref{eq:y-to-Phi}, the simplicity of the zeros of $\Phi_n$, and the argument used in the proof of Lemma~\ref{lem:log-bessel}, we obtain
 $z_j(\alpha)\rightarrow\rho_j$ as $\alpha\to\infty$. Since $z_j$ is strictly decreasing in $\alpha$ by \eqref{eq:zero-curve-derivative}, it follows that $z_j(\alpha)\downarrow\rho_j$ as $\alpha\to\infty$,
which proves \eqref{eq:sublinear-right-limit-general}.

We now retain the first perturbative term.  Define
\[
 y_\alpha:=\frac{w_\alpha-\Phi_n}{\varepsilon_\alpha}.
\]
Subtracting the equation for $\Phi_n$ from
\eqref{eq:sublinear-scaled} and dividing by
$\varepsilon_\alpha$ gives
\begin{equation}
 \begin{cases}
 y_\alpha''+\dfrac{n-1}{r}y_\alpha'+y_\alpha
 =|w_\alpha|^{q-2}w_\alpha,\qquad  r>0,\\[4pt]
 y_\alpha(0)=y_\alpha'(0)=0.
 \end{cases}
 \label{eq:y-alpha}
\end{equation}
The function $g(s)=|s|^{q-2}s$ is continuous, and hence uniformly continuous, on $[-1,1]$.
Consequently, \eqref{eq:upper1} and \eqref{eq:y-to-Phi} imply
$g(w_\alpha)\to g(\Phi_n)$ uniformly on every compact interval.
The Volterra formulation of \eqref{eq:y-alpha}, together with the
standard continuous-dependence estimate for the regular linear
equation, therefore gives
\begin{equation*}
 y_\alpha\rightarrow\Psi_n
 \quad\text{in }C^1_{\rm loc}([0,\infty)){\quad\text{as }\alpha\rightarrow\infty},
\end{equation*}
where $\Psi_n$ is the unique regular solution of
\begin{equation}
 \begin{cases}
 \Psi''+\dfrac{n-1}{r}\Psi'+\Psi
 =|\Phi_n|^{q-2}\Phi_n,\qquad  r>0,\\[4pt]
 \Psi(0)=\Psi'(0)=0.
 \end{cases}
 \label{eq:Psi-equation}
\end{equation}
For brevity, write $z_j=z_j(\alpha)$ and $\rho=\rho_j$.
{Since $w_\alpha(z_j)=0$,
\[
 \Phi_n(z_j)+\varepsilon_\alpha y_\alpha(z_j)=0.
\]
The local uniform boundedness of $y_\alpha$ gives
$|\Phi_n(z_j)|=O(\varepsilon_\alpha)$.  Because $\rho$ is a simple zero
and $z_j\to\rho$, this first implies
$z_j-\rho=O(\varepsilon_\alpha)$.  Applying the mean-value theorem and
using $y_\alpha(z_j)\to\Psi_n(\rho)$ then gives}
\begin{equation}
\frac{z_j-\rho}{\varepsilon_\alpha}
\longrightarrow
-\frac{\Psi_n(\rho)}{\Phi_n'(\rho)}\quad \text{as }\alpha\rightarrow\infty.
\label{eq:zero-shift-Psi}
\end{equation}

It remains to compute the constant on the right-hand side of \eqref{eq:zero-shift-Psi}. Multiply \eqref{eq:Psi-equation} by $\Phi_n$, multiply the equation for $\Phi_n$ by $\Psi_n$, and subtract. This gives %the Wronskian identity
\[
 \Bigl[r^{n-1}(\Phi_n\Psi_n'-\Phi_n'\Psi_n)\Bigr]'
 =r^{n-1}|\Phi_n|^q.
\]
{Integrating over $(0,\rho)$, the contribution at the origin
vanishes by regularity, and $\Phi_n(\rho)=0$. Hence}
\[
 -\rho^{n-1}\Phi_n'(\rho)\Psi_n(\rho)
 =\int_0^\rho r^{n-1}|\Phi_n(r)|^q\,dr.
\]
Since $\Phi_n'(\rho)\neq0$, it follows that
\[
 -\frac{\Psi_n(\rho)}{\Phi_n'(\rho)}
 =\frac{\displaystyle\int_0^\rho
 r^{n-1}|\Phi_n(r)|^q\,dr}
 {\rho^{n-1}\Phi_n'^2(\rho)}
 =\mathfrak c_{j,q,n}>0.
\]
Since $\varepsilon_\alpha=\alpha^{q-2}$,
\eqref{eq:zero-shift-Psi} is precisely
\eqref{eq:sublinear-first-order-zero}.
\end{proof}

\end{document}
